\documentclass[12pt,a4paper]{ctexart}

% ── Packages ──────────────────────────────────────────────────
\usepackage{amsmath,amssymb,amsthm}
\usepackage{bm}
\usepackage{geometry}
\usepackage{graphicx}
\usepackage{xcolor}
\usepackage{tikz}
\usetikzlibrary{positioning,arrows.meta}
\usepackage{enumitem}
\usepackage{hyperref}
\usepackage{tcolorbox}
\usepackage{fancyhdr}
\usepackage{booktabs}
\usepackage{float}
\usepackage{caption}
\usepackage{subcaption}
\usepackage{listings}
\usepackage{array}
\usepackage{multirow}

\geometry{margin=2.5cm}
\setlength{\headheight}{14.5pt}
\hypersetup{colorlinks=true,linkcolor=blue,citecolor=blue,urlcolor=blue,psdextra}

% ── Code style ────────────────────────────────────────────────
\lstset{
  basicstyle=\ttfamily\small,
  breaklines=true,
  frame=single,
  numbers=left,
  numberstyle=\tiny,
  keywordstyle=\color{blue},
  commentstyle=\color{gray},
  stringstyle=\color{red},
}

% ── Colored boxes ─────────────────────────────────────────────
\newtcolorbox{conceptbox}[1][]{
  colback=blue!5!white, colframe=blue!75!black,
  fonttitle=\bfseries, title=#1
}
\newtcolorbox{stepbox}[1][]{
  colback=orange!5!white, colframe=orange!75!black,
  fonttitle=\bfseries, title=#1
}
\newtcolorbox{warningbox}[1][]{
  colback=red!5!white, colframe=red!75!black,
  fonttitle=\bfseries, title=#1
}
\newtcolorbox{resultbox}[1][]{
  colback=green!5!white, colframe=green!50!black,
  fonttitle=\bfseries, title=#1
}

% ── Header / Footer ───────────────────────────────────────────
\pagestyle{fancy}
\fancyhf{}
\fancyhead[L]{\textbf{Nico\_Robo\_control --- 作业 7}}
\fancyhead[R]{2026}
\fancyfoot[C]{\thepage}

% ── Title ─────────────────────────────────────────────────────
\title{
  \textbf{作业 7：四自由度机械臂的计算力矩控制}\\[5pt]
  \large 逐步求解指南
}
\author{Nico\_Robo\_control}
\date{2026 年 6 月}

\begin{document}
\maketitle
\tableofcontents
\newpage

% ===================================================================
\section{问题概述}
% ===================================================================

我们考虑一个\ \textbf{四自由度}\ 机械臂，包含三个转动关节（关节 1--3）和一个移动关节（关节 4），与 \textbf{MLS 例 4.4}（第 177 页）中描述的完全一致。图~\ref{fig:manipulator}\ 展示了该机械臂在参考（零位）构型下的示意图。

\begin{figure}[H]
\centering
\begin{tikzpicture}[scale=1.2,>=stealth]
  % 基座
  \fill[gray!50] (-0.5,0) rectangle (0.5,-0.3);
  \draw[thick] (-0.5,0) -- (0.5,0);

  % 地面线
  \draw[thin] (-1,-0.3) -- (1,-0.3);
  \draw[thin] (-1,-0.3) -- (-1.2,-0.5);
  \draw[thin] (-0.5,-0.3) -- (-0.7,-0.5);
  \draw[thin] (0,-0.3) -- (0,-0.5);
  \draw[thin] (0.5,-0.3) -- (0.3,-0.5);
  \draw[thin] (1,-0.3) -- (0.8,-0.5);

  % 关节 1 (转动副, 原点)
  \draw[fill=blue!30] (0,0) circle (0.15);
  \node at (0,0.35) {$\theta_1$};
  \draw[->,red,thick] (0.3,0) arc (0:60:0.3);

  % 连杆 1 (水平)
  \draw[thick,fill=gray!20] (0,0) rectangle (4,0.2);
  \node at (2,0.5) {连杆 1 ($l_1=0.8$)};

  % 关节 2 (转动副)
  \draw[fill=blue!30] (4,0.1) circle (0.15);
  \node at (4,0.45) {$\theta_2$};
  \draw[->,red,thick] (4.3,0.1) arc (0:45:0.3);

  % 连杆 2 (水平)
  \draw[thick,fill=gray!20] (4,0.1) rectangle (7,0.3);
  \node at (5.5,0.6) {连杆 2 ($l_2=0.6$)};

  % 关节 3 (转动副)
  \draw[fill=blue!30] (7,0.2) circle (0.15);
  \node at (7,0.55) {$\theta_3$};
  \draw[->,red,thick] (7.3,0.2) arc (0:60:0.3);

  % 连杆 3 + 关节 4 (移动副，竖直)
  \draw[thick,fill=gray!20] (6.9,0.35) rectangle (7.1,1.5);
  \node at (7.5,0.9) {$d_4$ (移动)};
  \draw[<->,thick] (7.3,0.35) -- (7.3,1.5);

  % 末端执行器
  \draw[fill=red!50] (6.8,1.5) rectangle (7.2,1.7);

  % 坐标系
  \draw[->,thick] (0,0) -- (0.6,0) node[right] {$x$};
  \draw[->,thick] (0,0) -- (0,0.6) node[above] {$y$};
  \node at (-0.2,-0.2) {$z$ (向外)};
  \node at (-0.5,0) {$\{S\}$};
\end{tikzpicture}
\caption{四自由度机械臂（3R + 1P）参考构型示意图。三个转动关节轴线均平行于 $z$ 轴；移动关节沿 $z$ 方向平移。}
\label{fig:manipulator}
\end{figure}

\subsection{作业任务}

\begin{enumerate}[label=\textbf{(\alph*)}, leftmargin=*]
  \item 构建 \textbf{SIMULINK 计算力矩控制器}模型，跟踪\ \textbf{圆角阶跃函数}\ 形式的期望关节轨迹 $\bm{\theta}_d(t)$。期望终值：
  \[
  \theta_{d1}\to 2\,\text{rad},\quad
  \theta_{d2}\to 1\,\text{rad},\quad
  \theta_{d3}\to 0.5\,\text{rad},\quad
  d_{d4}\to 0.08\,\text{m}.
  \]
  \item 研究计算力矩控制的\ \textbf{跟踪鲁棒性}：将\ \textbf{被控对象}\ 的惯性参数扰动 5\%--10\%，同时保持\ \textbf{控制器}\ 参数为标称值。
  \item 实现并评估 \textbf{PD 控制}\ 方案，与计算力矩控制进行对比。
  \item 研究\ \textbf{力矩饱和}\ 效应：限制 $|\bm{\tau}| \le [6000,\;2000,\;1000,\;100]^\top$\,Nm（或 N）。
\end{enumerate}

% ===================================================================
\section{运动学建模}
% ===================================================================

\subsection{关节旋量（运动螺旋）}

采用\ \textbf{指数积（PoE）}\ 方法，每个关节用空间坐标系 $\{S\}$ 下的\ \textbf{旋量（screw）} $\bm{\xi}_i \in \mathbb{R}^6$ 表示。

对于轴线为 $\bm{\omega}$、轴上一点为 $\mathbf{q}$ 的\ \textbf{转动副}：
\[
\bm{\xi}_i = \begin{bmatrix} -\hat{\bm{\omega}}_i\,\mathbf{q}_i \\[2pt] \bm{\omega}_i \end{bmatrix}
\qquad\text{其中}\qquad
\text{hat}(\mathbf{v}) = \begin{bmatrix}
0 & -v_3 & v_2 \\
v_3 & 0 & -v_1 \\
-v_2 & v_1 & 0
\end{bmatrix}.
\]

对于平移方向为 $\mathbf{v}$ 的\ \textbf{移动副}：
\[
\bm{\xi}_i = \begin{bmatrix} \mathbf{v}_i \\[2pt] \mathbf{0}_{3\times 1} \end{bmatrix}.
\]

\begin{conceptbox}[关节旋量 $\bm{\xi}_i$]
\begin{align*}
\text{关节 1 (R, 绕 $z$ 轴, 过原点):}\quad
  \bm{\omega}_1 &= \begin{bmatrix}0\\0\\1\end{bmatrix},\;
  \mathbf{q}_1 = \begin{bmatrix}0\\0\\0\end{bmatrix}
  \;\Rightarrow\; \bm{\xi}_1 = \begin{bmatrix}0\\0\\0\\0\\0\\1\end{bmatrix} \\[5pt]
\text{关节 2 (R, 绕 $z$ 轴, 过 $[0,l_1,0]^\top$):}\quad
  \bm{\omega}_2 &= \begin{bmatrix}0\\0\\1\end{bmatrix},\;
  \mathbf{q}_2 = \begin{bmatrix}0\\l_1\\0\end{bmatrix}
  \;\Rightarrow\; \bm{\xi}_2 = \begin{bmatrix}l_1\\0\\0\\0\\0\\1\end{bmatrix} \\[5pt]
\text{关节 3 (R, 绕 $z$ 轴, 过 $[0,l_1+l_2,0]^\top$):}\quad
  \bm{\omega}_3 &= \begin{bmatrix}0\\0\\1\end{bmatrix},\;
  \mathbf{q}_3 = \begin{bmatrix}0\\l_1+l_2\\0\end{bmatrix}
  \;\Rightarrow\; \bm{\xi}_3 = \begin{bmatrix}l_1+l_2\\0\\0\\0\\0\\1\end{bmatrix} \\[5pt]
\text{关节 4 (P, 沿 $z$ 轴):}\quad
  \mathbf{v}_4 &= \begin{bmatrix}0\\0\\1\end{bmatrix}
  \;\Rightarrow\; \bm{\xi}_4 = \begin{bmatrix}0\\0\\1\\0\\0\\0\end{bmatrix}
\end{align*}
代入数值 $l_1 = 0.8$\,m, $l_2 = 0.6$\,m：
\[
\bm{\xi}_1 = \begin{bmatrix}0\\0\\0\\0\\0\\1\end{bmatrix},\quad
\bm{\xi}_2 = \begin{bmatrix}0.8\\0\\0\\0\\0\\1\end{bmatrix},\quad
\bm{\xi}_3 = \begin{bmatrix}1.4\\0\\0\\0\\0\\1\end{bmatrix},\quad
\bm{\xi}_4 = \begin{bmatrix}0\\0\\1\\0\\0\\0\end{bmatrix}.
\]
\end{conceptbox}

\subsection{指数积正运动学}

从空间坐标系到连杆 $i$ 的齐次变换为：
\[
\mathbf{g}_{sl_i}(\bm{\theta}) = e^{\hat{\bm{\xi}}_1\theta_1}\,e^{\hat{\bm{\xi}}_2\theta_2}\cdots e^{\hat{\bm{\xi}}_i\theta_i}\;\mathbf{g}_{sl_i}(0)
\]
其中 $\mathbf{g}_{sl_i}(0)$ 是连杆 $i$ 在参考构型下的位姿，$\hat{\bm{\xi}} \in \mathfrak{se}(3)$ 是旋量 $\bm{\xi}$ 的 $4\times4$ 矩阵表示：
\[
\hat{\bm{\xi}} = \begin{bmatrix} \hat{\bm{\omega}} & \mathbf{v} \\ \mathbf{0} & 0 \end{bmatrix}.
\]

\begin{stepbox}[参考位形 $\mathbf{g}_{sl_i}(0)$]
各连杆在零位时的参考坐标系为纯平移（无旋转）：
\begin{align*}
\mathbf{g}_{sl_1}(0) &= \begin{bmatrix} \mathbf{I} & [0,\; r_1,\; 0.3]^\top \\ \mathbf{0} & 1 \end{bmatrix}, \quad r_1 = l_1/2 = 0.4 \\[3pt]
\mathbf{g}_{sl_2}(0) &= \begin{bmatrix} \mathbf{I} & [0,\; l_1+r_2,\; 0.4]^\top \\ \mathbf{0} & 1 \end{bmatrix}, \quad r_2 = l_2/2 = 0.3 \\[3pt]
\mathbf{g}_{sl_3}(0) &= \begin{bmatrix} \mathbf{I} & [0,\; l_1+l_2,\; 0.5]^\top \\ \mathbf{0} & 1 \end{bmatrix} \\[3pt]
\mathbf{g}_{sl_4}(0) &= \begin{bmatrix} \mathbf{I} & [0,\; l_1+l_2,\; l_0]^\top \\ \mathbf{0} & 1 \end{bmatrix}, \quad l_0 = 0.2
\end{align*}
\end{stepbox}

\subsection{6×6 伴随变换}

伴随映射 $\mathrm{Ad}_{\mathbf{g}} \in \mathbb{R}^{6\times 6}$ 将旋量在不同坐标系间变换。对于 $\mathbf{g}=[\mathbf{R},\,\mathbf{p};\;\mathbf{0},\,1]$：
\[
\mathrm{Ad}_{\mathbf{g}} = \begin{bmatrix}
\mathbf{R} & \hat{\mathbf{p}}\mathbf{R} \\
\mathbf{0} & \mathbf{R}
\end{bmatrix}
\]
该公式在代码中由 \texttt{adjmg.m} 和 \texttt{MCNsim.m}（第 160--172 行）实现。

% ===================================================================
\section{动力学建模}
% ===================================================================

\subsection{连杆惯性参数}

\begin{table}[H]
\centering
\caption{各连杆质量和转动惯量}
\begin{tabular}{cccccc}
\toprule
连杆 & 质量 (kg) & $I_{xx}$ (kg·m$^2$) & $I_{yy}$ (kg·m$^2$) & $I_{zz}$ (kg·m$^2$) \\
\midrule
1 & $m_1 = 50$   & $m_1 l_1^2/12 = 2.667$ & $0.05 I_{x1}=0.133$ & $I_{x1}=2.667$ \\
2 & $m_2 = 30$   & $m_2 l_2^2/12 = 0.900$ & $0.05 I_{x2}=0.045$ & $I_{x2}=0.900$ \\
3 & $m_3 = 20$   & $2 I_{z3} = 0.180$     & $2 I_{z3} = 0.180$   & $0.1 I_{z2}=0.090$ \\
4 & $m_4 = 5$    & $2 I_{z4} = 0.108$     & $2 I_{z4} = 0.108$   & $0.6 I_{z3}=0.054$ \\
\bottomrule
\end{tabular}
\end{table}

连杆 $i$ 在其自身连体坐标系下的\ \textbf{广义惯性矩阵}\ 为 $6\times 6$ 矩阵：
\[
\hat{\mathbf{M}}_i = \begin{bmatrix}
m_i\mathbf{I}_3 & \mathbf{0} \\
\mathbf{0} & \mathbf{I}_i
\end{bmatrix},\qquad \mathbf{I}_i = \mathrm{diag}(I_{xi},\,I_{yi},\,I_{zi}).
\]

\subsection{空间惯性矩阵}

每个连体惯性矩阵通过伴随变换转换到空间坐标系 $\{S\}$：
\[
\hat{\mathbf{M}}_i^s = \mathrm{Ad}_{\mathbf{g}_{sl_i}(0)}^{-\top}\;\hat{\mathbf{M}}_i\;
\mathrm{Ad}_{\mathbf{g}_{sl_i}(0)}^{-1}.
\]
结果存储在三维数组 \texttt{calMp(:,:,i)} 中。

\subsection{运动方程}

机械臂动力学的标准形式：
\[
\boxed{\mathbf{M}(\bm{\theta})\,\ddot{\bm{\theta}} + \mathbf{C}(\bm{\theta},\dot{\bm{\theta}})\,\dot{\bm{\theta}} + \mathbf{N}(\bm{\theta},\dot{\bm{\theta}}) = \bm{\tau}}
\]
其中：
\begin{itemize}
  \item $\mathbf{M}(\bm{\theta}) \in \mathbb{R}^{4\times 4}$ --- \textbf{质量/惯性矩阵}（对称正定）
  \item $\mathbf{C}(\bm{\theta},\dot{\bm{\theta}}) \in \mathbb{R}^{4\times 4}$ --- \textbf{科氏力/离心力矩阵}
  \item $\mathbf{N}(\bm{\theta},\dot{\bm{\theta}}) \in \mathbb{R}^{4}$ --- \textbf{重力 + 摩擦力向量}
  \item $\bm{\tau} \in \mathbb{R}^{4}$ --- \textbf{关节力矩/力}
\end{itemize}

\subsection{\texorpdfstring{旋量法计算 $\mathbf{M}(\bm{\theta})$}{旋量法计算 M(theta)}}

采用\ \textbf{Christoffel 符号}\ 方法和空间惯性矩阵形式：
\[
M_{ij}(\bm{\theta}) = \sum_{l = \max(i,j)}^{n} \bm{\xi}_i^\top\,\mathbf{A}_{li}^\top\,
\hat{\mathbf{M}}_l^s\,\mathbf{A}_{lj}\,\bm{\xi}_j
\]
其中 $\mathbf{A}_{ij} \in \mathbb{R}^{6\times 6}$ 是关节间的伴随映射：
\[
\mathbf{A}_{ij} = \begin{cases}
\mathrm{Ad}^{-1}_{\mathbf{g}_{ji}}, & i > j \\[3pt]
\mathbf{I}_6, & i = j \\[3pt]
\mathbf{0}, & i < j
\end{cases},\qquad
\mathbf{g}_{ji} = \mathbf{g}_j^{-1}\,\mathbf{g}_i.
\]

\subsection{\texorpdfstring{计算科氏矩阵 $\mathbf{C}(\bm{\theta},\dot{\bm{\theta}})$}{计算科氏矩阵 C(theta, dtheta)}}

科氏矩阵通过 $\mathbf{M}$ 的偏导数利用\ \textbf{第一类 Christoffel 符号}\ 求得：
\[
C_{ij} = \sum_{k=1}^{n} \Gamma_{ijk}\,\dot{\theta}_k,\qquad
\Gamma_{ijk} = \frac{1}{2}\left(
\frac{\partial M_{ij}}{\partial\theta_k} +
\frac{\partial M_{ik}}{\partial\theta_j} -
\frac{\partial M_{kj}}{\partial\theta_i}
\right).
\]

偏导数 $\partial M_{ij}/\partial\theta_k$ 的解析计算公式为：
\begin{align*}
\frac{\partial M_{ij}}{\partial\theta_k} = \sum_{l=\max(i,j)}^{n}
\Bigl[
&\bigl(\mathrm{ad}_{\mathbf{A}_{ki}\bm{\xi}_i}\,\bm{\xi}_k\bigr)^\top\,
\mathbf{A}_{lk}^\top\,\hat{\mathbf{M}}_l^s\,\mathbf{A}_{lj}\,\bm{\xi}_j \\[3pt]
+\;&\bm{\xi}_i^\top\,\mathbf{A}_{li}^\top\,\hat{\mathbf{M}}_l^s\,
\mathbf{A}_{lk}\,\bigl(\mathrm{ad}_{\mathbf{A}_{kj}\bm{\xi}_j}\,\bm{\xi}_k\bigr)
\Bigr],
\end{align*}
该式在 $k \le l$ 时成立，否则为零。其中 $\mathrm{ad}_{\bm{\xi}}$ 是 $6\times 6$ 的李括号（伴随）矩阵：
\[
\mathrm{ad}_{\bm{\xi}} = \begin{bmatrix}
\hat{\bm{\omega}} & \hat{\mathbf{v}} \\
\mathbf{0} & \hat{\bm{\omega}}
\end{bmatrix}.
\]

\subsection{\texorpdfstring{计算 $\mathbf{N}(\bm{\theta},\dot{\bm{\theta}})$}{计算 N(theta, dtheta)}}

$\mathbf{N}$ 向量由两部分组成：

\subsubsection*{1. 重力项}
重力势能的梯度：
\[
\frac{\partial V}{\partial\theta_i} = -\sum_{l=i}^{n} \bm{\xi}_i^\top\,
\mathrm{Ad}_{\mathbf{g}_{0l}}^{-\top}\,\hat{\mathbf{M}}_l^s\,
\mathbf{G}_0
\]
其中 $\mathbf{G}_0 \in \mathbb{R}^6$ 是空间坐标系下的\ \textbf{重力旋量}：
\[
\mathbf{G}_0 = \begin{bmatrix} -\mathbf{g}_{\text{accel}} \\ \mathbf{0} \end{bmatrix}
= \begin{bmatrix} 0\\0\\9.81\\0\\0\\0 \end{bmatrix}.
\]
（符号遵循代码约定：线速度部分取 $-\mathbf{g}_{\text{accel}}$。）

\subsubsection*{2. 黏性摩擦项}
\[
N_{\text{fric},i} = \beta_i\,\dot{\theta}_i
\]
其中 $\bm{\beta} = [1.0,\;0.5,\;0.4,\;1.0]$（转动关节单位为 Nm·s/rad，移动关节为 N·s/m）。

\subsubsection*{总计：}
\[
\mathbf{N}(\bm{\theta},\dot{\bm{\theta}}) = -\nabla_\theta V(\bm{\theta}) + \mathrm{diag}(\bm{\beta})\,\dot{\bm{\theta}}.
\]

\subsection{代码实现}

动力学计算在 \texttt{MCN.m} 中实现（其自包含版本 \texttt{MCNsim.m} 用于 Simulink）。\texttt{MCNsim.m} 文件在 Simulink MATLAB Function 模块中使用，因为它将 \texttt{adjmg} 作为局部嵌套函数（无外部依赖）。

\begin{lstlisting}[language=Matlab, caption={MCNsim.m --- 动力学计算（用于 Simulink 的自包含版本）}]
function [Ms, Cs, Ns] = MCNsim(xi, calMp, lM, beta, m4g, theta, dtheta)
    n = length(theta);
    % 步骤 1：正运动学 --- 计算所有 g_{0i}
    g = cell(1, n+1); g{1} = eye(4);
    for i = 1:n
        g{i+1} = g{i} * expm(hat_se3(xi(:,i)) * theta(i));
    end
    % 步骤 2：伴随映射 A_{ij}
    A = cell(n,n);
    for i = 1:n, for j = 1:n
        if i > j,      A{i,j} = inv(adjmg(inv(g{j+1}) * g{i+1}));
        elseif i == j, A{i,j} = eye(6);
        else,          A{i,j} = zeros(6);
        end
    end, end
    % 步骤 3：质量矩阵 M
    for i = 1:n, for j = 1:n
        M_ij = 0;
        for l = lM(i,j):n
            M_ij = M_ij + xi(:,i)' * A{l,i}' * calMp(:,:,l) * A{l,j} * xi(:,j);
        end
        Ms(i,j) = M_ij;
    end, end
    % 步骤 4：通过 Christoffel 符号计算科氏矩阵 C
    for i = 1:n, for j = 1:n
        C_ij = 0;
        for k = 1:n
            C_ij = C_ij + 0.5 * (dM(i,j,k) + dM(i,k,j) - dM(k,j,i)) * dtheta(k);
        end
        Cs(i,j) = C_ij;
    end, end
    % 步骤 5：计算 N（重力 + 摩擦）
    grad_V = compute_gravity_potential_gradient(xi, calMp, g, n);
    Ns = -grad_V + beta(:) .* dtheta(:);
end
\end{lstlisting}

\subsection{\texorpdfstring{耦合矩阵 $l_M$}{耦合矩阵 lM}}

矩阵 \texttt{lM} 编码了哪些连杆对 $(i,j)$ 对质量矩阵有贡献：
\[
l_M(i,j) = \max(i,j).
\]
对于 $n=4$：
\[
l_M = \begin{bmatrix}
1 & 2 & 3 & 4 \\
2 & 2 & 3 & 4 \\
3 & 3 & 3 & 4 \\
4 & 4 & 4 & 4
\end{bmatrix}.
\]

% ===================================================================
\section{第 (a) 部分：计算力矩控制器}
% ===================================================================

\subsection{控制律}

\textbf{计算力矩}\ （反馈线性化）控制器施加：
\[
\boxed{\bm{\tau} = \mathbf{M}(\bm{\theta})\,\ddot{\bm{\theta}}^* + \mathbf{C}(\bm{\theta},\dot{\bm{\theta}})\,\dot{\bm{\theta}} + \mathbf{N}(\bm{\theta},\dot{\bm{\theta}})}
\]
其中 $\ddot{\bm{\theta}}^*$ 是为实现跟踪而选择的虚拟输入。代入动力学方程可得\ \textbf{双积分}\ 误差动力学：
\[
\ddot{\bm{\theta}} = \ddot{\bm{\theta}}^*.
\]

$\ddot{\bm{\theta}}^*$ 的\ \textbf{PD + 前馈}\ 控制律为：
\[
\boxed{\ddot{\bm{\theta}}^* = \ddot{\bm{\theta}}_d + \mathbf{K}_v\bigl(\dot{\bm{\theta}}_d - \dot{\bm{\theta}}\bigr) + \mathbf{K}_p\bigl(\bm{\theta}_d - \bm{\theta}\bigr)}
\]
其中 $\mathbf{K}_p = \mathrm{diag}(k_{p1},\ldots,k_{p4})$ 和 $\mathbf{K}_v = \mathrm{diag}(k_{v1},\ldots,k_{v4})$ 为正定增益矩阵。

由此可得\ \textbf{线性误差动力学}：
\[
\ddot{\mathbf{e}} + \mathbf{K}_v\,\dot{\mathbf{e}} + \mathbf{K}_p\,\mathbf{e} = \mathbf{0},
\qquad \mathbf{e}(t) = \bm{\theta}_d(t) - \bm{\theta}(t).
\]

\subsection{增益选取}

对于期望的闭环固有频率 $\omega_n$ 和阻尼比 $\zeta$：
\[
k_{pi} = \omega_{n,i}^2,\qquad k_{vi} = 2\zeta_i\omega_{n,i}.
\]

建议的初始值：
\[
\mathbf{K}_p = \mathrm{diag}(2500,\;2500,\;2500,\;1600),\qquad
\mathbf{K}_v = \mathrm{diag}(100,\;100,\;100,\;80)
\]
（对应转动关节 $\omega_n \approx 50$\,rad/s，移动关节约 40\,rad/s，$\zeta=1$）。

\subsection{圆角阶跃轨迹}

由于计算力矩控制需要 $\bm{\theta}_d$ 为\ \textbf{二次可微}\ 函数，这里使用\ \textbf{五次多项式}\ 圆角阶跃：
\[
r(t) = \begin{cases}
0, & t \le t_0 \\[3pt]
10\left(\dfrac{t-t_0}{\Delta t}\right)^3 - 15\left(\dfrac{t-t_0}{\Delta t}\right)^4 + 6\left(\dfrac{t-t_0}{\Delta t}\right)^5, & t_0 < t < t_0+\Delta t \\[3pt]
1, & t \ge t_0 + \Delta t
\end{cases}
\]

\begin{conceptbox}[五次圆角阶跃函数的性质]
\begin{itemize}
  \item $r(t_0) = 0$， $r(t_0+\Delta t) = 1$
  \item $\dot{r}(t_0) = \dot{r}(t_0+\Delta t) = 0$
  \item $\ddot{r}(t_0) = \ddot{r}(t_0+\Delta t) = 0$
  \item 该函数是处处 $\mathcal{C}^2$ 连续的。
\end{itemize}
\end{conceptbox}

过渡参数（来自 \texttt{initialize\_simulink.m}）：
\[
t_0 = 1.0\,\text{s},\qquad \Delta t = 0.2\,\text{s},
\]
即阶跃从 $t=1.0$\,s 开始，在 $t=1.0$ 到 $t=1.2$\,s 之间过渡，然后保持恒定。

期望关节轨迹通过缩放得到：
\[
\bm{\theta}_d(t) = \begin{bmatrix} 2 \\ 1 \\ 0.5 \\ 0.08 \end{bmatrix} \cdot r(t),\qquad
\dot{\bm{\theta}}_d(t) = \begin{bmatrix} 2 \\ 1 \\ 0.5 \\ 0.08 \end{bmatrix} \cdot \dot{r}(t),\qquad
\ddot{\bm{\theta}}_d(t) = \begin{bmatrix} 2 \\ 1 \\ 0.5 \\ 0.08 \end{bmatrix} \cdot \ddot{r}(t).
\]

\subsection{Simulink 实现}

\begin{stepbox}[Simulink 逐步搭建指南]
\begin{enumerate}[leftmargin=*, label=\textbf{步骤 \arabic*}., ref=步骤 \arabic*]
  \item \textbf{复制圆角阶跃模块}：从 \texttt{MLS4\_4\_rounded\_step\_inputR2024a.slx} 复制到你的新模型中。该模块输出 $[y,\;\dot{y},\;\ddot{y}]^\top$。
  \item \textbf{缩放输出：} 将 $y$、$\dot{y}$、$\ddot{y}$ 信号分别乘以幅值向量 $[2,\,1,\,0.5,\,0.08]$，得到 $\bm{\theta}_d$、$\dot{\bm{\theta}}_d$、$\ddot{\bm{\theta}}_d$。

  \item \textbf{实现两个 MATLAB Function 模块}\ 用于动力学计算：
  \begin{itemize}
    \item \textbf{控制器动力学模块：} 包含使用\ \textbf{标称}\ 参数的 \texttt{MCNsim}（命名为 \texttt{MCN\_ctrl}）。输入：$\bm{\theta}$、$\dot{\bm{\theta}}$。输出：$\mathbf{M}_c$、$\mathbf{C}_c$、$\mathbf{N}_c$。
    \item \textbf{被控对象动力学模块：} 包含使用（可能扰动的）参数的 \texttt{MCNsim}（命名为 \texttt{MCN\_plant}）。输入：$\bm{\theta}$、$\dot{\bm{\theta}}$。输出：$\mathbf{M}_p$、$\mathbf{C}_p$、$\mathbf{N}_p$。
  \end{itemize}
  （初始时，两个模块使用完全相同的标称参数。）

  \item \textbf{计算 $\ddot{\bm{\theta}}^*$：}
  \[
  \ddot{\bm{\theta}}^* = \ddot{\bm{\theta}}_d + \mathbf{K}_v(\dot{\bm{\theta}}_d - \dot{\bm{\theta}}) + \mathbf{K}_p(\bm{\theta}_d - \bm{\theta})
  \]
  使用矩阵增益模块或 MATLAB Function 模块。

  \item \textbf{计算控制力矩 $\bm{\tau}$：}
  \[
  \bm{\tau} = \mathbf{M}_c\,\ddot{\bm{\theta}}^* + \mathbf{C}_c\,\dot{\bm{\theta}} + \mathbf{N}_c
  \]

  \item \textbf{求解正动力学}\ 使用被控对象模型：
  \[
  \ddot{\bm{\theta}} = \mathbf{M}_p^{-1}\bigl(\bm{\tau} - \mathbf{C}_p\,\dot{\bm{\theta}} - \mathbf{N}_p\bigr)
  \]
  使用 MATLAB Function 模块对 $\mathbf{M}_p$ 求逆并计算加速度。

  \item \textbf{积分两次}\ 得到 $\dot{\bm{\theta}}$ 和 $\bm{\theta}$：
  \begin{itemize}
    \item 第一个积分器：$\ddot{\bm{\theta}} \to \dot{\bm{\theta}}$
    \item 第二个积分器：$\dot{\bm{\theta}} \to \bm{\theta}$
  \end{itemize}

  \item \textbf{初始条件：} 设置 $\bm{\theta}(0) = [\pi/4,\;\pi/6,\;\pi/2.5,\;-0.04]^\top$，$\dot{\bm{\theta}}(0)=\mathbf{0}$。
\end{enumerate}
\end{stepbox}

\begin{figure}[H]
\centering
\begin{tikzpicture}[node distance=2.5cm,>=stealth, auto,
  block/.style={draw, rectangle, minimum width=2.5cm, minimum height=1cm, align=center, fill=blue!10},
  sum/.style={draw, circle, minimum size=0.7cm}]

  % 圆角阶跃
  \node [block] (step) {圆角阶跃\\$[y,\dot{y},\ddot{y}]^\top$};
  \node [block, right=2cm of step] (scale) {按幅值缩放\\$[2,1,0.5,0.08]$};

  % PD + 前馈
  \node [sum, right=2cm of scale] (sum1) {$+$};
  \node [sum, below=1cm of sum1] (sum2) {$+$};
  \node [sum, below=1cm of sum2] (sum3) {$+$};

  % 增益
  \node [block, right=1.5cm of sum1] (Kp) {$\mathbf{K}_p$};
  \node [block, right=1.5cm of sum2] (Kv) {$\mathbf{K}_v$};

  % 求和节点 tau*
  \node [sum, right=1.5cm of Kv] (sum_tau) {$+$};

  % 控制器 MCN
  \node [block, below=2cm of scale] (mcn_ctrl) {MCN\_ctrl\\(标称参数)};

  % tau 计算
  \node [block, right=1.5cm of sum_tau] (tau_block) {$\bm{\tau} = \mathbf{M}_c\ddot{\bm{\theta}}^*$\\$+\mathbf{C}_c\dot{\bm{\theta}}+\mathbf{N}_c$};

  % 被控对象正动力学
  \node [block, right=1.5cm of tau_block] (plant) {正动力学\\$\ddot{\bm{\theta}}=\mathbf{M}_p^{-1}(\bm{\tau}-\mathbf{C}_p\dot{\bm{\theta}}-\mathbf{N}_p)$};

  % 积分器
  \node [block, right=1.5cm of plant] (int2) {$\int$};
  \node [block, right=1.5cm of int2] (int1) {$\int$};

  \node [right=1.5cm of int1] (theta_out) {$\bm{\theta}$};
  \node [above=0.8cm of int2] (dtheta_out) {$\dot{\bm{\theta}}$};

  % 连线
  \draw[->] (step) -- (scale);
  \draw[->] (scale.east) -- node[above] {$\bm{\theta}_d$} (sum1);
  \draw[->] (scale.east) -- ++(0,-1) -- node[above] {$\dot{\bm{\theta}}_d$} (sum2);
  \draw[->] (scale.east) -- ++(0,-2) -- node[above] {$\ddot{\bm{\theta}}_d$} (sum3);

  \draw[->] (sum1) -- (Kp);
  \draw[->] (sum2) -- (Kv);
  \draw[->] (Kp) -- (sum_tau);
  \draw[->] (Kv) -- (sum_tau);
  \draw[->] (sum3) -- (sum_tau);

  % 反馈
  \draw[->] (int2) -- ++(0,3.5) -- node[above] {$\dot{\bm{\theta}}$} ++(-14,0) -- (sum2);
  \draw[->] (int1) -- ++(0,5) -- node[above] {$\bm{\theta}$} ++(-17,0) -- (sum1);

  % MCN 连接
  \draw[->] (mcn_ctrl) -- (tau_block);
  \draw[->] (sum_tau) -- (tau_block);
  \draw[->] (tau_block) -- (plant);
  \draw[->] (plant) -- (int2);
  \draw[->] (int2) -- node[above] {$\dot{\bm{\theta}}$} (int1);
  \draw[->] (int1) -- (theta_out);

\end{tikzpicture}
\caption{计算力矩控制器 Simulink 高层框图。使用两个 MCN 模块：一个用于控制器（标称参数），另一个用于被控对象（可引入扰动）。}
\label{fig:simulink_block}
\end{figure}

\subsection{验证}

\begin{stepbox}[验证步骤]
\begin{enumerate}
  \item 运行仿真 3--5 秒。
  \item 检查所有关节是否收敛到期望终值：$\theta_1\to 2$, $\theta_2\to 1$, $\theta_3\to 0.5$, $d_4\to 0.08$。
  \item 验证跟踪误差 $\mathbf{e}(t) = \bm{\theta}_d(t) - \bm{\theta}(t)$ 是否收敛到零。
  \item 检查力矩 $\bm{\tau}$ 是否在合理范围内。
\end{enumerate}
\end{stepbox}

% ===================================================================
\section{第 (b) 部分：鲁棒性研究}
% ===================================================================

\subsection{方法}

为评估鲁棒性，我们在\ \textbf{被控对象}\ （实际物理机器人）和\ \textbf{控制器}\ （假定标称参数）之间引入\ \textbf{不匹配}。

\begin{conceptbox}[扰动策略]
\textbf{仅被控对象：} 在 MCN\_plant 模块中，将每个连杆的\ \textbf{空间惯性矩阵 $\hat{\mathbf{M}}_i^s$}\ 的所有元素乘以 $(1+\delta)$，其中 $\delta \in \{-0.10,\,-0.05,\,+0.05,\,+0.10\}$。

\textbf{控制器：} MCN\_ctrl 模块保持使用原始（标称）$\hat{\mathbf{M}}_i^s$。
\end{conceptbox}

\subsection{Simulink 实现方法}

在被控对象使用的 \texttt{MCNsim} 函数中，修改 \texttt{calMp} 数组：
\begin{lstlisting}[language=Matlab]
% 仅在 MCN_plant 中：
delta = 0.05;  % 或 -0.05, 0.10, -0.10
calMp_perturbed = calMp * (1 + delta);
% 后续动力学计算中使用 calMp_perturbed
\end{lstlisting}

或者，将 \texttt{delta} 参数作为额外输入传入你的 MATLAB Function 模块。

\subsection{需要观察和报告的内容}

\begin{enumerate}[leftmargin=*]
  \item \textbf{跟踪误差：} 绘制 $\delta=0$（基线）与 $\delta=\pm 5\%$、$\pm 10\%$ 下的 $\mathbf{e}(t)$ 曲线。参数不匹配时，误差\ \textbf{不应收敛到零}\ --- 这正是计算力矩控制的关键弱点（依赖精确的模型对消）。
  \item \textbf{稳态误差：} 即使到达稳态，由于重力补偿不正确，非零误差仍然存在。
  \item \textbf{瞬态行为：} 超调量和调节时间随参数扰动变化。
  \item \textbf{影响方向：} 正 $\delta$（机器人比模型重）和负 $\delta$（机器人比模型轻）会产生不同的误差模式。
\end{enumerate}

\begin{resultbox}[预期结果]
\begin{itemize}
  \item \textbf{无扰动}\ ($\delta = 0$) 时：完美渐近跟踪，$\mathbf{e}(t) \to \mathbf{0}$。
  \item \textbf{有扰动}\ ($\delta \neq 0$) 时：非零稳态误差；$|\delta|$ 越大，误差越大。
  \item 移动关节（关节 4）由于直接重力负载，通常最为敏感。
\end{itemize}
\end{resultbox}

% ===================================================================
\section{第 (c) 部分：PD 控制方案}
% ===================================================================

\subsection{PD 控制律}

\textbf{PD + 重力补偿}\ 控制器为：
\[
\boxed{\bm{\tau}_{\text{PD}} = \mathbf{K}_p\bigl(\bm{\theta}_d - \bm{\theta}\bigr) + \mathbf{K}_v\bigl(\dot{\bm{\theta}}_d - \dot{\bm{\theta}}\bigr) + \mathbf{N}(\bm{\theta})}
\]

与计算力矩控制相比，PD 控制器：
\begin{itemize}
  \item \textbf{不使用} $\mathbf{M}(\bm{\theta})$ 和 $\mathbf{C}(\bm{\theta},\dot{\bm{\theta}})$（仅用 $\mathbf{N}$ 做重力补偿）。
  \item 实现\ \textbf{更简单}，计算量更小。
  \item 闭环误差动力学是\ \textbf{非线性}\ 的（无法完美对消非线性项）。
  \item \textbf{不需要} $\ddot{\bm{\theta}}_d$（只需 $\bm{\theta}_d$ 和 $\dot{\bm{\theta}}_d$）。
\end{itemize}

PD 控制下的闭环动力学：
\[
\mathbf{M}(\bm{\theta})\,\ddot{\bm{\theta}} + \mathbf{C}(\bm{\theta},\dot{\bm{\theta}})\,\dot{\bm{\theta}} + \mathbf{N}(\bm{\theta}) = \mathbf{K}_p\mathbf{e} + \mathbf{K}_v\dot{\mathbf{e}} + \mathbf{N}(\bm{\theta})
\]
该方程\ \textbf{不能}\ 简化为简单的线性系统。

\subsection{Simulink 实现}

\begin{stepbox}[PD 控制器框图 --- 比计算力矩更简单]
\begin{enumerate}
  \item 移除 $\mathbf{M}_c$、$\mathbf{C}_c$ 和 $\mathbf{N}_c$ 模块。
  \item 移除前馈 $\ddot{\bm{\theta}}_d$ 项。
  \item 计算：
  \[
  \bm{\tau} = \mathbf{K}_p(\bm{\theta}_d - \bm{\theta}) + \mathbf{K}_v(\dot{\bm{\theta}}_d - \dot{\bm{\theta}}) + \mathbf{N}_p(\bm{\theta})
  \]
  其中 $\mathbf{N}_p$ 仅为重力向量（摩擦可选）。注意：若仅用 PD（无重力补偿），则省略 $\mathbf{N}_p$ 项。
  \item 将 $\bm{\tau}$ 直接输入正动力学（被控对象模型）。
\end{enumerate}
\end{stepbox}

\subsection{PD 增益整定}

PD 增益通常需要比计算力矩增益\ \textbf{更大}，以补偿模型解耦的缺失：
\[
\mathbf{K}_p^{\text{PD}} \approx (2\text{--}5) \times \mathbf{K}_p^{\text{CT}},\qquad
\mathbf{K}_v^{\text{PD}} \approx (2\text{--}5) \times \mathbf{K}_v^{\text{CT}}.
\]

\subsection{对比内容与报告}

\begin{table}[H]
\centering
\caption{对比：计算力矩控制 vs.\ PD 控制}
\begin{tabular}{p{3.5cm}p{6cm}p{5cm}}
\toprule
\textbf{对比项} & \textbf{计算力矩} & \textbf{PD（+ 重力补偿）} \\
\midrule
模型依赖 & 完整 $\mathbf{M},\mathbf{C},\mathbf{N}$ & 仅 $\mathbf{N}$（或全无） \\
计算开销 & 高（需矩阵求逆） & 低 \\
跟踪精度 & 优秀（模型精确时） & 中等（存在非线性耦合） \\
参数误差鲁棒性 & 差（依赖精确对消） & 好（无需精确对消） \\
增益整定 & 系统化（极点配置） & 启发式 / 试凑法 \\
需要 $\ddot{\bm{\theta}}_d$？ & 是 & 否 \\
\bottomrule
\end{tabular}
\end{table}

\begin{resultbox}[预期结果]
\begin{itemize}
  \item \textbf{计算力矩}\ 在模型精确时实现更快、更准确的跟踪。
  \item \textbf{PD}\ 在瞬态过程中因未补偿的动态耦合而具有较大的跟踪误差。
  \item \textbf{PD}\ 对参数不确定性更鲁棒（误差退化是渐进的，而非突变）。
  \item 有重力补偿时，PD 对阶跃参考可实现零稳态误差（由无源性理论保证）。
  \item 无重力补偿时，PD 存在稳态误差。
\end{itemize}
\end{resultbox}

% ===================================================================
\section{第 (d) 部分：力矩饱和}
% ===================================================================

\subsection{问题描述}

实际执行器具有有限的力矩/力输出能力。饱和限值为：
\[
|\tau_i| \le \tau_{i,\max},\qquad
\bm{\tau}_{\max} = \begin{bmatrix} 6000 \\ 2000 \\ 1000 \\ 100 \end{bmatrix}
\;\text{Nm（关节 1--3）/ N（关节 4）}.
\]

当计算力矩控制器要求的 $\tau_i$ 超过 $\tau_{i,\max}$ 时，实际施加的力矩会被\ \textbf{截断}：
\[
\tau_i^{\text{实际}} = \mathrm{sat}(\tau_i^{\text{指令}};\;-\tau_{i,\max},\;\tau_{i,\max})
= \begin{cases}
\tau_{i,\max}, & \tau_i^{\text{指令}} > \tau_{i,\max} \\
\tau_i^{\text{指令}}, & |\tau_i^{\text{指令}}| \le \tau_{i,\max} \\
-\tau_{i,\max}, & \tau_i^{\text{指令}} < -\tau_{i,\max}
\end{cases}
\]

\subsection{Simulink 实现}

\begin{stepbox}[饱和模块的放置]
在计算力矩输出与正动力学被控对象输入之间插入一个\ \textbf{Saturation 模块}\ （位于 Simulink $\to$ Discontinuities 库），限值设为 $\pm[6000,\,2000,\,1000,\,100]$。
\end{stepbox}

\subsection{饱和效应}

\begin{conceptbox}[关键现象]
\begin{enumerate}
  \item \textbf{积分器 Windup（卷绕）：} 当力矩饱和时，误差继续累积在积分器（如使用）中，导致退出饱和时产生大幅超调。这称为\ \textbf{windup 现象}。
  \item \textbf{响应变慢：} 系统无法按规定加速度运动，瞬态过程中的跟踪性能下降。
  \item \textbf{潜在不稳定：} 如果控制器依赖非线性项的精确对消，而实际施加的力矩与指令力矩差异显著，误差动力学不再是线性的，可能变得不稳定。
  \item \textbf{阶跃幅度依赖性：} 更大的期望阶跃需要更大的力矩；饱和限值可能被更严重地突破。
\end{enumerate}
\end{conceptbox}

\subsection{抗 Windup 策略（选做研究）}

\begin{enumerate}[leftmargin=*]
  \item \textbf{反算（Back-calculation）：} 将差值 $(\bm{\tau}_{\text{实际}} - \bm{\tau}_{\text{指令}})$ 反馈以减小积分作用。
  \item \textbf{条件积分：} 在饱和激活时冻结积分器。
  \item \textbf{参考指令调节器：} 当指令力矩超过限值时，缩小 $\ddot{\bm{\theta}}^*$。
\end{enumerate}

\subsection{需要报告的内容}

\begin{enumerate}[leftmargin=*]
  \item 绘制各关节的指令力矩 vs.\ 实际力矩。
  \item 对比有饱和和无饱和时的跟踪性能。
  \item 识别哪些关节最先饱和（通常是支撑整个手臂的关节 1）。
  \item 研究饱和限值是否足以完成期望轨迹。
  \item 如观察到 windup 现象，建议并测试抗 windup 改进方法。
\end{enumerate}

\begin{resultbox}[预期结果]
\begin{itemize}
  \item 关节 1（基座关节）需要最大力矩，最可能首先饱和。
  \item 饱和时，跟踪误差增大且可能不收敛到零。
  \item 响应变慢，可能因 windup 出现超调。
  \item 加入抗 windup 后，超调减小，但速度仍受执行器能力限制。
\end{itemize}
\end{resultbox}

% ===================================================================
\section{完整 Simulink 模型总览}
% ===================================================================

\begin{figure}[H]
\centering
\fbox{%
\begin{minipage}{0.9\textwidth}
\vspace{5pt}
\textbf{控制器侧（每个时间步执行）：}
\begin{enumerate}[leftmargin=*, nosep]
  \item 读取 $\bm{\theta}$、$\dot{\bm{\theta}}$（来自积分器）
  \item 读取 $\bm{\theta}_d$、$\dot{\bm{\theta}}_d$、$\ddot{\bm{\theta}}_d$（来自圆角阶跃生成器）
  \item 计算 $\mathbf{e} = \bm{\theta}_d - \bm{\theta}$， $\dot{\mathbf{e}} = \dot{\bm{\theta}}_d - \dot{\bm{\theta}}$
  \item 计算：$\ddot{\bm{\theta}}^* = \ddot{\bm{\theta}}_d + \mathbf{K}_v\dot{\mathbf{e}} + \mathbf{K}_p\mathbf{e}$
  \item 调用 MCN\_ctrl$(\bm{\theta},\dot{\bm{\theta}})$ $\to$ $\mathbf{M}_c,\mathbf{C}_c,\mathbf{N}_c$
  \item 计算：$\bm{\tau} = \mathbf{M}_c\ddot{\bm{\theta}}^* + \mathbf{C}_c\dot{\bm{\theta}} + \mathbf{N}_c$
  \item 施加饱和：$\bm{\tau}_a = \mathrm{sat}(\bm{\tau})$
\end{enumerate}
\vspace{3pt}
\hrule
\vspace{3pt}
\textbf{被控对象侧（正动力学）：}
\begin{enumerate}[leftmargin=*, nosep]
  \item 调用 MCN\_plant$(\bm{\theta},\dot{\bm{\theta}})$ $\to$ $\mathbf{M}_p,\mathbf{C}_p,\mathbf{N}_p$
  \item 计算：$\ddot{\bm{\theta}} = \mathbf{M}_p^{-1}(\bm{\tau}_a - \mathbf{C}_p\dot{\bm{\theta}} - \mathbf{N}_p)$
  \item 积分：$\dot{\bm{\theta}} = \int \ddot{\bm{\theta}}\,dt$， $\bm{\theta} = \int \dot{\bm{\theta}}\,dt$
\end{enumerate}
\vspace{5pt}
\end{minipage}}
\caption{控制+仿真系统的完整计算流程。}
\label{fig:complete_flow}
\end{figure}

% ===================================================================
\section{期望绘制的图表}
% ===================================================================

报告中需包含以下图形：

\begin{enumerate}[leftmargin=*]
  \item \textbf{跟踪性能（标称情况）：} 全部 4 个关节的 $\bm{\theta}(t)$ vs.\ $\bm{\theta}_d(t)$，以及跟踪误差 $\mathbf{e}(t)$。
  \item \textbf{控制力矩（标称）：} 全部 4 个关节的 $\bm{\tau}(t)$。
  \item \textbf{鲁棒性：} $\delta = -10\%, -5\%, 0, +5\%, +10\%$ 的跟踪误差（叠加或子图）。
  \item \textbf{PD 对比：} PD 与计算力矩的跟踪性能对比（叠加图）。
  \item \textbf{力矩饱和：} 指令力矩 vs.\ 实际力矩，以及有/无饱和时的跟踪误差。
  \item （选做）\textbf{抗 Windup：} 饱和条件下有/无抗 windup 的跟踪误差对比。
\end{enumerate}

% ===================================================================
\section{核心公式速查}
% ===================================================================

\begin{conceptbox}[计算力矩控制]
\[
\bm{\tau} = \mathbf{M}(\bm{\theta})\bigl[\ddot{\bm{\theta}}_d + \mathbf{K}_v\dot{\mathbf{e}} + \mathbf{K}_p\mathbf{e}\bigr] + \mathbf{C}(\bm{\theta},\dot{\bm{\theta}})\dot{\bm{\theta}} + \mathbf{N}(\bm{\theta},\dot{\bm{\theta}})
\]
\[
\mathbf{e} = \bm{\theta}_d - \bm{\theta},\qquad \dot{\mathbf{e}} = \dot{\bm{\theta}}_d - \dot{\bm{\theta}}
\]
\end{conceptbox}

\begin{conceptbox}[PD 控制 + 重力补偿]
\[
\bm{\tau} = \mathbf{K}_p\mathbf{e} + \mathbf{K}_v\dot{\mathbf{e}} + \mathbf{N}_{\text{grav}}(\bm{\theta})
\]
\end{conceptbox}

\begin{conceptbox}[正动力学（仿真用）]
\[
\ddot{\bm{\theta}} = \mathbf{M}^{-1}(\bm{\theta})\bigl[\bm{\tau} - \mathbf{C}(\bm{\theta},\dot{\bm{\theta}})\dot{\bm{\theta}} - \mathbf{N}(\bm{\theta},\dot{\bm{\theta}})\bigr]
\]
\end{conceptbox}

\begin{conceptbox}[五次圆角阶跃]
\[
r(t) = 10\left(\frac{t-t_0}{\Delta t}\right)^3 - 15\left(\frac{t-t_0}{\Delta t}\right)^4 + 6\left(\frac{t-t_0}{\Delta t}\right)^5
\]
对于 $t_0 < t < t_0+\Delta t$，其中 $t_0=1.0$\,s，$\Delta t=0.2$\,s。
\end{conceptbox}

% ===================================================================
\section{常见问题排查}
% ===================================================================

\begin{warningbox}[常见陷阱]
\begin{enumerate}[leftmargin=*]
  \item \textbf{代数环（Algebraic Loop）：} 如果 $\mathbf{M}$、$\mathbf{C}$ 或 $\mathbf{N}$ 依赖于 $\ddot{\bm{\theta}}$，Simulink 将报代数环错误。确保你的 MCN 函数仅以 $(\bm{\theta},\dot{\bm{\theta}})$ 为输入（不包含 $\ddot{\bm{\theta}}$）。
  \item \textbf{维度不匹配：} 所有信号均为 $4\times 1$ 向量。请小心使用 Selector 模块、Demux/Mux 模块。
  \item \textbf{$\mathbf{M}$ 奇异：} 该机械臂的 $\mathbf{M}(\bm{\theta})$ 始终可逆。如果出现奇异矩阵警告，请检查惯性参数。
  \item \textbf{MATLAB Function 模块依赖：} 在 Simulink MATLAB Function 模块内使用\ \texttt{MCNsim.m}\ （自包含版本），而非 \texttt{MCN.m}（有外部依赖）。
  \item \textbf{步长过大：} 如果求解器发散，在 Configuration Parameters 中减小最大步长（建议尝试 1e-3 或更小）。
  \item \textbf{初始条件：} 初始 $\bm{\theta}$ 必须在积分器模块中设置。初始 $\dot{\bm{\theta}}$ 通常为零。
  \item \textbf{控制器/被控对象混淆：} 第 (b) 部分中，仅在 MCN\_plant（被控对象）中扰动参数，MCN\_ctrl（控制器）保持原值不变。
\end{enumerate}
\end{warningbox}

% ===================================================================
\section{总结}
% ===================================================================

本次作业涵盖了机器人控制系统设计的完整流程：

\begin{enumerate}[leftmargin=*]
  \item \textbf{运动学与动力学建模}\ 采用旋量理论（指数积）方法，提供了一种优雅的、不依赖坐标系的公式化方法，适合符号和数值计算。
  \item \textbf{计算力矩控制}\ 在模型精确时通过完全对消非线性动力学并施加线性误差动力学，实现完美跟踪。
  \item \textbf{鲁棒性}\ 是计算力矩控制的致命弱点 --- 参数不匹配导致非零跟踪误差，这促使后续引入自适应或鲁棒控制扩展。
  \item \textbf{PD 控制}\ 更简单、更鲁棒，但因未补偿的动态耦合而牺牲了跟踪性能。
  \item \textbf{执行器饱和}\ 施加了必须考虑的实践限制，可能需要抗 windup 策略。
\end{enumerate}

\vspace{1cm}
\begin{center}
\rule{0.5\textwidth}{0.5pt}\\[5pt]
\textit{求解指南结束}
\end{center}

\end{document}
